(1) 介质本构 \(\boldsymbol P=\varepsilon_0(\varepsilon_r-1)\boldsymbol E_{\rm in}\)，而内部总场 \(\boldsymbol E_{\rm in}=\boldsymbol E_0-\boldsymbol L\boldsymbol P/\varepsilon_0\)。逐主轴消元得 \(\boxed{E_{{\rm in},i}=E_{0,i}/[1+L_i(\varepsilon_r-1)]}\) 与 \(P_i=\varepsilon_0(\varepsilon_r-1)E_{0,i}/[1+L_i(\varepsilon_r-1)]\)。体积 \(V=4\pi abc/3\)，均匀极化总偶极矩 \(p_i=VP_i=\alpha_iE_{0,i}\)，故 \(\boxed{\alpha_i=\varepsilon_0V(\varepsilon_r-1)/[1+L_i(\varepsilon_r-1)]}\)。形状造成张量极化率，即使材料本身各向同性。

(2) 球形时积分化为 \(L_i=(R^3/2)\int_0^\infty(R^2+s)^{-5/2}ds=(R^3/2)\cdot(2/(3R^3))=1/3\)。代入 \(\alpha_i\) 得 \(\alpha=4\pi\varepsilon_0R^3(\varepsilon_r-1)/(\varepsilon_r+2)\)，与静电介电球边界匹配一致。若 \(\varepsilon_r\to\infty\)，各主轴 \(\boxed{\alpha_i\to\varepsilon_0V/L_i}\)；球形导体恢复 \(4\pi\varepsilon_0R^3\)。极化率有限来自形状退极化，而非认为导体内部场仍有限。

(3) 线性诱导偶极在固定外场下可逆能量为 \(U=-\frac12\boldsymbol E_0\cdot\boldsymbol\alpha\boldsymbol E_0\)，因从零场逐渐极化需要对场积分，不能用永久偶极的 \(-pE\) 少掉系数 \(1/2\)。故 \(U(\theta)=-E_0^2[\alpha_x\cos^2\theta+\alpha_y\sin^2\theta]/2\)。转矩 \(\boxed{\tau_z=-dU/d\theta=-(E_0^2/2)(\alpha_x-\alpha_y)\sin2\theta}\)。若 \(\alpha_x>\alpha_y\)，\(\theta=0,\pi\) 为能量最低的稳定对齐，\(\theta=\pi/2\) 为不稳定方向。

(4) 长旋转椭球给定公式在 \(e\to0\) 用 \(\ln[(1+e)/(1-e)]=2(e+e^3/3+e^5/5+\cdots)\)，得 \(L_x=1/3-2e^2/15+O(e^4)\)。当 \(e\to1^-\)，\((1-e^2)\ln[2/(1-e)]\to0\)，故 \(L_x\to0\)，而 \(L_y=L_z=(1-L_x)/2\to1/2\)。因此长轴极化率 \(\alpha_x\) 比横轴更大，导体极限尤其明显。极端尖端会出现强局域场，有限击穿、非线性极化、材料尺寸效应均使理想均匀线性椭球近似失效。
